\documentclass[12pt]{article}
\usepackage[utf8]{inputenc}
\usepackage[T1]{fontenc}
\usepackage{hyperref}
\title{Amam Language Resources for K-12 Curriculum Development}
\author{Compiled by Claude Sonnet 5.5 (claude-sonnet-5-5) for Andr\'as Kornai}
\date{}
\begin{document}
\maketitle

\section{Introduction}
Amam (ISO 639-3 \texttt{wer}; Glottocode weri1253) is a Goilalan (Kunimaipan) language of Morobe Province, Papua New Guinea. The ISO code \texttt{wer} nominally covers Weri and Amam. Search results indicate that a 2013 grammar by Aki and Pennington found the two to be separate languages, and Glottolog likewise lists Amam as its own entry. Very little material exists. Nothing found in search shows Amam used in schooling.

\section{Primary Educational Resources}
I found no primers, textbooks or school materials for Amam in web search.

\section{Linguistic and Reference Materials}
\begin{itemize}
\item Kunimaipan languages, Wikipedia overview of the family: \url{https://en.wikipedia.org/wiki/Kunimaipan_languages}
\item Glottolog entry, Spoken L1 Language: Amam: \url{https://glottolog.org/resource/languoid/id/amam1246}
\item ASJP word-list entry for Amam: \url{https://asjp.clld.org/languages/AMAM}
\end{itemize}
Search summaries mention a 2013 tentative grammar by Aki and Pennington, based on spoken and written texts, with a corpus of about 2500 words and a noun class system of at least 15 classes. I did not see a direct, verified URL for the grammar itself, so it is not listed with a link.

\section{Religious and Community Resources}
I found nothing for this section in search results.

\section{Digital and Audio Resources}
I found no audio, video or software resources specific to Amam.

\section{Curriculum Development Context}
Documentation is thin, and I saw no school programme. A curriculum would first need community-led orthography and materials development, drawing on the Aki and Pennington grammar and any recordings behind it. Tok Pisin and English would be the likely bridging languages.
\end{document}